\section{Re-cognition: The Closing and Reopening of the Turn}
\label{sec:recognition}

\noindent
The cycle turns to re-cognition, the moment at which the revised and reopened appraisal carried out of reflection becomes the cognition with which a further turn begins. The task of this moment is the least like a discrete operation and the most like a passage: it is the carrying forward of what the turn has accumulated into the cognition of the next, and it is the point at which the lifecycle is seen to be a spiral rather than a circle. The moment is brief to state and consequential in what it establishes about the cycle as a whole.

\subsection{Re-cognition is not a return to the initial cognition}
\label{sec:rec-spiral}

The cognition with which a turn of the cycle ends is not the cognition with which it began. The turn has passed through practice, through the evaluation of what the practice generated, and through the reflection that read the evaluation against the value it could not reach; and the cognition that emerges carries the accumulated difference. What the parties take their flourishing to consist in, at the close of a turn, has been altered by the turn: deepened where the turn ascended, narrowed where it declined, and in either case not identical to the cognition that opened it. Re-cognition is this altered cognition, and the lifecycle is therefore not a circle that returns to its starting point but a spiral that returns to the starting configuration bearing the phase the turn accumulated.

\begin{claim}[Re-cognition carries the holonomy of the turn]
The cognition with which a turn of the cycle closes differs from the cognition with which it opened by the phase the turn has accumulated, and re-cognition is the carrying of this difference into the next turn. The lifecycle is accordingly a spiral and not a circle: it returns to the moment of cognition, but to a cognition altered by the turn, and the accumulated alteration over many turns is the long ascent or the long decline of a shared life. A turn that ascends returns a deepened cognition of the good, enriched by what the turn generated and corrected by what reflection found; a turn that declines returns a narrowed cognition, impoverished by what the turn depleted; and a turn traversed without ascent or decline returns the cognition unchanged, the closed circle of a repetition that goes round without rising. Re-cognition is the moment at which the phase of the turn is registered in the cognition of the next, and through which the holonomy of the lifecycle accumulates across the history of a relation.
\end{claim}

\subsection{The accumulation across the time of a relation}
\label{sec:rec-accumulation}

The registration of each turn's phase in the cognition of the next is the mechanism by which a shared life rises or declines across its history. A relation is many nested turns (\xref{sec:cyc-time}), and the re-cognition that closes each turn sets the cognition of the next, so that the phases compound: a relation whose turns severally ascend compounds its ascent, its cognition of the good deepening turn upon turn, its practice enriching, its parties developing through the shared life into more than they were; a relation whose turns severally decline compounds its decline, however intact its outward form. The passage from courtship through marriage into the later years is, on this account, the compounding of the phases of its turns, and the difference between a shared life that has flourished across its arc and one that has declined across it is the accumulated holonomy of its many turns, registered at each re-cognition and carried into each successive cognition.

This is the point at which the lifecycle's relation to the time of a real shared life is most concrete. The marriage at which a relation arrives is not the completion of the cycle but a re-cognition of unusual depth, a turn at which the cognition of the good is reconceived under the altered conditions of a life now formally and practically joined; and the later years are the compounding of the turns that follow it. The cognition of flourishing is, across this arc, never fixed and never finally attained, but reconceived at each re-cognition in the light of what the prior turns have made of the parties and of their shared life. The good of a shared life is, in this respect, not a destination the cycle approaches but the very substance of its turning, generated and reconceived at each turn and accumulated, for better or worse, across them all.

\subsection{The openness of the spiral}
\label{sec:rec-openness}

The spiral does not close. Re-cognition opens a further turn, and there is no final re-cognition at which the cognition of the good is fixed once and for all and the cycle comes to rest in a completed understanding. To suppose such a final re-cognition would be to suppose a state of achieved and possessed flourishing, the closed circle of a cognition that no longer revises, which is, by the argument of \xref{sec:cyc-persistence}, not the completion of the relation but its cessation. The good of a shared life is unguaranteeable and processual, in the sense the prior paper established, and its cognition is correspondingly unfinished: reconceived at each turn, never possessed as a final understanding, and carried forward into a further turn that the cycle, so long as the relation lives, continues to open. Re-cognition is therefore not the cycle's conclusion but its renewal, and the lifecycle, considered as a whole, is an open spiral that the relation sustains by continuing to turn. The paper now considers this whole, in the three regimes its turning may assume.
